\chapter{Módulos sobre un dominio de ideales
principales}\label{ch:b3:modules}

Álgebra lineal sobre un \emph{anillo} en lugar de sobre un cuerpo: este
pequeño cambio de hipótesis produce uno de los grandes teoremas de
unificación del álgebra. Un \omterm{def:b3:modules:module}{módulo} sobre
$\Z$ es un grupo abeliano; un \omterm{def:b3:modules:module}{módulo}
sobre $K[X]$ es un espacio vectorial dotado de un endomorfismo. El
teorema de estructura de los \omterm{def:b3:modules:free}{módulos
finitamente generados} sobre un \omterm{def:b3:rings:pidufd}{DIP}
clasifica, pues, de un solo golpe, todos los grupos abelianos
\omterm{def:b3:modules:free}{finitamente generados} \emph{y} todos los
endomorfismos salvo semejanza --- la reducción de Jordan, que en segundo
año se obtuvo mediante delicadas inducciones, cae como corolario, junto
con su hermana más sutil, la forma canónica \emph{racional}, válida
sobre todo cuerpo. El motor de cálculo es la
\omterm{thm:b3:modules:smith}{forma normal de Smith}, una aritmética de
matrices digna de Euclides.

En todo el capítulo, $A$ es un anillo conmutativo, pronto un
\omterm{def:b3:rings:pidufd}{DIP};
«\omterm{def:b3:modules:module}{módulo}» significa $A$-módulo.

\section{Módulos y módulos libres}

\begin{definition}\label{def:b3:modules:module}
Un \emph{$A$-módulo}\index{módulo} es un grupo abeliano $(M, +)$ con una
multiplicación por escalares $A \times M \to M$ que cumple los axiomas
de espacio vectorial: $a(x + y) = ax + ay$, $(a + b)x = ax + bx$,
$(ab)x = a(bx)$, $1x = x$. Los submódulos, los cocientes $M/N$, los
morfismos (aplicaciones $A$-lineales), las sumas directas
$\bigoplus_i M_i$ y los teoremas de isomorfía se definen y se demuestran
palabra por palabra igual que para los espacios vectoriales y los grupos
abelianos; en particular, $M/\ker f \cong
\operatorname{im} f$ para todo morfismo $f$.
\end{definition}

\begin{example}\label{ex:b3:modules:examples}
Los tres casos que motivan la teoría.
\begin{enumerate}
  \item $A = K$ un cuerpo: los \omterm{def:b3:modules:module}{módulos}
  son los espacios vectoriales. \item $A = \Z$: los
  \omterm{def:b3:modules:module}{módulos} son exactamente los grupos
  abelianos ($nx$ ha de ser $x + \dots + x$) y los submódulos son los
  subgrupos. \item $A = K[X]$: un \omterm{def:b3:modules:module}{módulo}
  es un $K$-espacio vectorial $V$ junto con la aplicación $K$-lineal
  $u\colon x \mapsto X\cdot x$ --- recíprocamente, todo par $(V, u)$ con
  $u \in \mathcal L(V)$ se convierte en un $K[X]$-módulo mediante
  $P \cdot x = P(u)(x)$. Los submódulos son precisamente los subespacios
  $u$-estables.
\end{enumerate}
Un \omterm{def:b3:rings:ideal}{ideal} de $A$ es exactamente un submódulo
de $A$; un \omterm{def:b3:rings:ideal}{anillo cociente} $A/I$ es un
$A$-módulo. A diferencia de los espacios vectoriales, los
\omterm{def:b3:modules:module}{módulos} pueden tener
\emph{\omterm{def:b3:modules:torsion}{torsión}}: en $\Z/6\Z$, el
elemento $\bar 3 \ne 0$ es aniquilado por $2 \neq 0$.
\end{example}

\begin{definition}\label{def:b3:modules:free}
$M$ es \emph{finitamente generado}\index{módulo finitamente generado} si
$M = Ax_1 + \dots + Ax_n$ para ciertos $x_i$. $M$ es
\emph{libre}\index{módulo libre} de rango $n$ si $M \cong A^n$, es
decir, si tiene una \emph{base} (una familia generadora $A$-linealmente
independiente). Todo $M$ finitamente generado es cociente de un
\omterm{def:b3:modules:module}{módulo} libre:
$(a_1, \dots, a_n) \mapsto \sum a_ix_i$ aplica $A^n$ sobre $M$.
\end{definition}

\begin{proposition}[Invariancia del rango]\label{prop:b3:modules:rank}
Si $A \neq 0$ y $A^m \cong A^n$, entonces $m = n$.
\end{proposition}

\begin{proof}
Tómese un \omterm{def:b3:rings:primemaximal}{ideal maximal}
$\mathfrak m$ de $A$ (el~\cref{thm:b3:rings:krull}) y póngase
$k = A/\mathfrak m$, que es un cuerpo. Un isomorfismo $f \colon A^m
\to A^n$ lleva $\mathfrak m A^m$ en $\mathfrak m A^n$ (linealidad),
luego induce un isomorfismo de cocientes
\[
A^m/\mathfrak m A^m \;\cong\; A^n/\mathfrak m A^n,
\qquad\text{es decir}\qquad k^m \cong k^n
\]
como $k$-espacios vectoriales (el cociente $A^m/\mathfrak m A^m$ es
aniquilado por $\mathfrak m$, de modo que la acción de $A$ se factoriza
a través de $k$; las imágenes de la base canónica forman una $k$-base).
La teoría de la dimensión sobre el cuerpo $k$ da $m = n$.
\end{proof}

\begin{theorem}[Submódulos de los módulos
libres]\label{thm:b3:modules:submodule}
Sea $A$ un \omterm{def:b3:rings:pidufd}{DIP} y sea $M \subseteq A^n$ un
submódulo. Entonces $M$ es \omterm{def:b3:modules:free}{libre} de rango
$\leq n$.
\end{theorem}

\begin{proof}
Inducción sobre $n$. Para $n = 1$: $M$ es un
\omterm{def:b3:rings:ideal}{ideal}, luego $M = (0)$
(\omterm{def:b3:modules:free}{libre} de rango $0$) o $M = dA \cong A$
($x \mapsto dx$ es inyectiva: dominio). Para $n > 1$: sea
$\pi \colon A^n \to A$ la última coordenada. Entonces $\pi(M)$ es un
\omterm{def:b3:rings:ideal}{ideal}, $(0)$ o $dA$. Si $(0)$:
$M \subseteq A^{n-1} \times \{0\}$ y se aplica la hipótesis de
inducción. En caso contrario, elíjase $x_0 \in M$ con $\pi(x_0) = d$.
Todo $x \in M$ se escribe de manera única
\[
x = \underbrace{\frac{\pi(x)}{d}}_{\in A}\, x_0 +
\Bigl(x - \tfrac{\pi(x)}d x_0\Bigr),
\qquad x - \tfrac{\pi(x)}d x_0 \in M \cap \ker\pi
\]
($\pi(x) \in dA$, luego el coeficiente está en $A$). Así $M = Ax_0
\oplus (M \cap \ker \pi)$: la suma es directa, pues $\pi(ax_0) =
ad = 0$ obliga a $a = 0$. Por inducción, $M \cap \ker\pi \subseteq
\ker \pi \cong A^{n-1}$ es \omterm{def:b3:modules:free}{libre} de rango
$\leq n - 1$; añadiendo $x_0$ (independiente de $\ker\pi$, como acabamos
de ver) se obtiene una base de $M$ de cardinal $\leq n$.
\end{proof}

\begin{remark}\label{rem:b3:modules:presentation}
En consecuencia, sobre un \omterm{def:b3:rings:pidufd}{DIP} todo
\omterm{def:b3:modules:free}{módulo finitamente generado} $M$ admite una
\emph{presentación finita}: una sobreyección $\varphi\colon A^n \to M$
tiene núcleo \omterm{def:b3:modules:free}{libre} con base
$c_1, \dots, c_k$ ($k \leq n$), y $M
\cong A^n / \operatorname{im}(C)$, donde $C \in M_{n,k}(A)$ es la matriz
cuyas columnas son los $c_j$. Entender $M$ significa entender una
\emph{matriz sobre $A$} salvo cambios de base en el origen y en el
destino --- el asunto de la sección siguiente.
\end{remark}

\section{Forma normal de Smith}

\begin{definition}\label{def:b3:modules:equivalence}
Dos matrices $B, C \in M_{n,k}(A)$ son \emph{equivalentes} si $C =
QBP$ con $Q \in GL_n(A)$, $P \in GL_k(A)$ (invertibles \emph{sobre $A$}:
determinante en $A^\times$). Matrices de presentación equivalentes
definen \omterm{def:b3:modules:module}{módulos} $A^n/
\operatorname{im}$ isomorfos (cámbiense las bases de $A^n$ y $A^k$).
\end{definition}

\begin{theorem}[Forma normal de Smith]\label{thm:b3:modules:smith}
Sea $A$ un \omterm{def:b3:rings:pidufd}{DIP} y sea $B \in M_{n,k}(A)$.
Entonces $B$ es equivalente a una matriz diagonal\index{forma normal de
Smith}
\[
\operatorname{diag}(d_1, d_2, \dots, d_r, 0, \dots, 0),
\qquad d_1 \mid d_2 \mid \cdots \mid d_r \neq 0 ,
\]
y los $d_i$ son únicos salvo
\omterm{def:b3:rings:divisibility}{asociados}: $d_1 \cdots d_i$ es un
máximo común divisor de los menores $i \times i$ de $B$ (en particular,
ese máximo común divisor es un invariante de la equivalencia). Los $d_i$
son los \emph{factores invariantes}\index{factores invariantes} de $B$.
\end{theorem}

\begin{proof}
\emph{Existencia.} Si $B = 0$, hemos terminado. En caso contrario,
considérese el conjunto de los \omterm{def:b3:rings:ideal}{ideales}
$(b)$ generados por las entradas de las matrices equivalentes a $B$;
como $A$ es \omterm{def:b3:rings:noetherian}{noetheriano}, elíjanse una
matriz $B'$ equivalente a $B$ y una entrada $d$ de $B'$ con $(d)$
\emph{maximal} en ese conjunto. Llévese $d$ a la posición $(1,1)$
mediante intercambios de filas y columnas.

\emph{Afirmación: $d$ divide a toda entrada de $B'$.} Primero, la
columna 1: si $b_{i1}$ no es múltiplo de $d$, sea
$e = \gcd(d, b_{i1}) = ud +
vb_{i1}$ (Bézout), de modo que $(e) \supsetneq (d)$. El truco de la
matriz $2 \times 2$: actuando sobre las filas $1$ y $i$ por
\[
\begin{pmatrix} u & v \\ -\dfrac{b_{i1}}e & \dfrac de
\end{pmatrix}
\in GL_2(A)
\qquad
\Bigl(\det = \tfrac{ud + v b_{i1}}e = 1\Bigr)
\]
se obtiene una matriz equivalente con entrada $e$ en la posición
$(1,1)$, lo que contradice la maximalidad de $(d)$. Así que $d$ divide
la columna $1$ y, simétricamente, la fila $1$. Restando múltiplos de la
fila 1 y de la columna 1 se anulan: $B'$ es equivalente a
$\begin{pmatrix} d &
0\\ 0 & B''\end{pmatrix}$. A continuación, $d$ divide toda entrada $b$
de $B''$: súmese la fila de $b$ a la fila 1 (una operación elemental; la
nueva primera fila contiene $d$ y entradas $b$) y repítase el argumento
de anulación de columnas: un no múltiplo mejoraría de nuevo $(d)$. Ahora
inducción sobre el tamaño: $B''$, cuyas entradas son todas divisibles
por $d$, tiene una forma de Smith $\operatorname{diag}(d_2,
\dots)$ cuyas entradas siguen siendo divisibles por $d$ (toda entrada de
cualquier $QB''P$ es una $A$-combinación de entradas de $B''$); póngase
$d_1 =
d$.

\emph{Unicidad.} Denotemos por $D_i(B)$ un máximo común divisor de todos
los menores $i \times i$. Las operaciones de filas y columnas y, más en
general, la multiplicación por \emph{cualquier} matriz, no pueden
reducir ese máximo común divisor: los menores $i \times i$ de $QB$ son
$A$-combinaciones de los de $B$ (desarrollo de Cauchy--Binet; o
directamente: cada fila de $QB$ es combinación de filas de $B$, y los
menores son multilineales en las filas). Así que $D_i(QBP)$ y $D_i(B)$
se dividen mutuamente: $D_i$ es un invariante de la equivalencia. Sobre
la forma diagonal, los menores $i
\times i$ no nulos son los productos de $i$ de los $d_j$, y la
divisibilidad $d_1 \mid \dots \mid d_r$ hace de $d_1 \cdots d_i$ el
máximo común divisor. Por tanto $d_1 \cdots d_i = D_i(B)$ salvo
unidades, y $d_i =
D_i/D_{i-1}$ queda determinado.
\end{proof}

\begin{method}\label{met:b3:modules:smith}
Sobre un \omterm{def:b3:rings:pidufd}{dominio euclídeo} ($\Z$, $K[X]$),
la reducción de Smith es un algoritmo --- no hace falta el argumento de
maximalidad: llévese a la posición $(1,1)$ la entrada de menor tamaño
euclídeo; si no divide alguna entrada de su fila o de su columna, una
división euclídea deja allí un resto estrictamente menor ---
intercámbiese y vuélvase a empezar (termina porque los tamaños
decrecen); cuando divide toda su fila y su columna, anúlense estas; si
no divide alguna entrada interior, súmese esa fila a la fila $1$ y
vuélvase a empezar; recúrrase sobre el bloque interior. En la práctica,
con matrices enteras: calcúlense $D_1 = \gcd$ de las entradas, $D_2$,
\dots\ mediante menores para tamaños pequeños, o ejecútese el algoritmo.
\end{method}

\begin{example}[Una reducción de Smith, con todo detalle]
\label{ex:b3:modules:smithworked}
Reduzcamos $M = \begin{pmatrix} 1 & 2 & 3\\ 4 & 5 & 6\\ 7 & 8 & 9
\end{pmatrix}$ sobre $\Z$. La esquina $1$ lo divide todo: anúlense su
fila y su columna ($L_2 \leftarrow L_2 - 4L_1$, $L_3
\leftarrow L_3 - 7L_1$, y después $C_2 \leftarrow C_2 - 2C_1$, $C_3
\leftarrow C_3 - 3C_1$):
\[
M \sim \begin{pmatrix} 1 & 0 & 0\\ 0 & -3 & -6\\ 0 & -6 & -12
\end{pmatrix} .
\]
En el bloque interior, la esquina $-3$ divide todas las entradas: $L_3
\leftarrow L_3 - 2L_2$ y $C_3 \leftarrow C_3 - 2C_2$ lo dejan reducido a
$\operatorname{diag}(-3, 0)$. Ajustando los signos (multiplíquese una
fila por $-1$, operación lícita):
\[
M \sim \operatorname{diag}(1, 3, 0),
\qquad
\Z^3/M\Z^3 \cong \Z/3\Z \times \Z .
\]
Comprobación mediante los divisores determinantales: $D_1 =
\gcd(\text{entradas}) = 1$; todo menor $2\times2$ de $M$ es múltiplo de
$3$ (por ejemplo, $\det\bigl(\begin{smallmatrix}1 & 2\\ 4
& 5\end{smallmatrix}\bigr) = -3$) y uno de ellos vale $-3$: $D_2 =
3$; $D_3 = \det M = 0$. Por tanto $d_1 = 1$, $d_2 = 3$, $d_3 = 0$: la
misma respuesta. Dos lecciones: un factor invariante nulo registra la
caída del rango (el conúcleo gana un sumando
\omterm{def:b3:modules:free}{libre} $\Z$), y la cadena de divisibilidad
$1 \mid 3 \mid 0$ es el certificado de Smith --- una reducción diagonal
que viola la cadena (por ejemplo $\operatorname{diag}(2, 3)$, que un
descuido puede producir a partir de
$\bigl(\begin{smallmatrix}2 & 0\\ 0 &
3\end{smallmatrix}\bigr)$ deteniéndose demasiado pronto: la forma de
Smith correcta es $\operatorname{diag}(1, 6)$, ¡pues aquí $D_1 = 1$!) no
está terminada.
\end{example}

\begin{omfigure}
\begin{tikzpicture}[scale=0.72]
  \clip (-3.4,-1.3) rectangle (4.4,4.3);
  \foreach \x in {-3,...,4}
    \foreach \y in {-1,...,4}
      \fill[black!40] (\x,\y) circle (1.6pt);
  \fill[omDef!12, opacity=0.7] (0,0) -- (1,3) -- (3,3) -- (2,0)
    -- cycle;
  \foreach \a/\b in {0/0, 2/0, 4/0, -2/0, 1/3, 3/3, -1/3, -3/3,
                     2/6, -2/-6, 0/6}
    \fill[omThm] (\a,\b) circle (2.6pt);
  \draw[omThm, thick, ->] (0,0) -- (2,0)
    node[midway, below, font=\small] {$(2,0)$};
  \draw[omThm, thick, ->] (0,0) -- (1,3)
    node[midway, left, font=\small] {$(1,3)$};
\end{tikzpicture}

{\small El subretículo $L = \Z(2,0) + \Z(1,3)$ de $\Z^2$ (puntos rojos).
La \omterm{thm:b3:modules:smith}{forma normal de Smith} de
$\bigl(\begin{smallmatrix} 2 &
1\\ 0 & 3\end{smallmatrix}\bigr)$ es $\operatorname{diag}(1, 6)$: en la
base adaptada $f_1 = (1,3)$, $f_2 = (0,1)$ de $\Z^2$ se tiene
$L = \Z f_1 \oplus \Z\,6f_2$, luego $\Z^2/L \cong \Z/6\Z$ --- el índice
es igual a $\abs{\det} = 6$, el área del dominio fundamental sombreado.}
\end{omfigure}

\section{El teorema de estructura}

\begin{definition}\label{def:b3:modules:torsion}
Sea $A$ un dominio y sea $M$ un $A$-módulo. El submódulo de
\emph{torsión}\index{torsión} es
\[
T(M) = \{x \in M : ax = 0 \text{ para cierto } a \neq 0\}
\]
(es un submódulo: si $ax = by = 0$, entonces $ab(x + y) = 0$,
$ab \ne 0$). $M$ es \emph{sin torsión} si
$T(M) = 0$, y un \emph{\omterm{def:b3:modules:module}{módulo} de
torsión} si $T(M) = M$.
\end{definition}

\begin{theorem}[Estructura de los módulos finitamente generados sobre un
DIP]\label{thm:b3:modules:structure} Sea $A$ un
\omterm{def:b3:rings:pidufd}{DIP} y sea $M$ un $A$-módulo
\omterm{def:b3:modules:free}{finitamente generado}. Existen un
$r \in \N$ único y elementos $d_1 \mid d_2 \mid
\cdots \mid d_s$ no nulos y no unidades, únicos salvo
\omterm{def:b3:rings:divisibility}{asociados}, tales que\index{teorema
de estructura de los módulos}
\[
M \;\cong\; A^r \,\oplus\, A/(d_1) \oplus \cdots \oplus A/(d_s).
\]
Además, $T(M) \cong A/(d_1)\oplus\dots\oplus A/(d_s)$ y $M/T(M)
\cong A^r$: un \omterm{def:b3:modules:free}{módulo
finitamente generado} y
\emph{sin \omterm{def:b3:modules:torsion}{torsión}} sobre un
\omterm{def:b3:rings:pidufd}{DIP} es
\omterm{def:b3:modules:free}{libre}.
\end{theorem}

\begin{proof}
\emph{Existencia.} Preséntese $M \cong A^n/\operatorname{im}(C)$
(la~\cref{rem:b3:modules:presentation}) y llévese $C$ a forma de Smith:
tras los dos cambios de base, $M \cong A^n / (d_1A \times \dots
\times d_rA \times 0 \times \dots \times 0) \cong A/(d_1) \oplus
\dots \oplus A/(d_r) \oplus A^{\,n - r}$. Descártense los factores en
los que $d_i$ es una unidad ($A/(d_i) = 0$); la cadena de divisibilidad
se conserva.

\emph{Identificación de la \omterm{def:b3:modules:torsion}{torsión}.} En
la descomposición, $A^r$ es sin \omterm{def:b3:modules:torsion}{torsión}
(un dominio no tiene divisores de cero) y cada $A/(d_i)$ es de
\omterm{def:b3:modules:torsion}{torsión} (aniquilado por $d_i \ne 0$);
una suma directa descompone la \omterm{def:b3:modules:torsion}{torsión}
en consecuencia: $T(M) = \bigoplus_i A/(d_i)$ y $M/T(M) \cong A^r$.

\emph{Unicidad de $r$}: $M/T(M) \cong A^r$ solo depende de $M$, y
la~\cref{prop:b3:modules:rank} lo determina $r$.

\emph{Unicidad de los $d_i$}: basta tratar el
\omterm{def:b3:modules:torsion}{módulo de torsión} $T = T(M)$.
Descompóngase cada $d_i$ en primos y sepárese mediante el teorema chino
del resto (el~\cref{thm:b3:rings:crt}; primos distintos generan
\omterm{thm:b3:rings:crt}{ideales comaximales}):
\[
A/(d) \cong \bigoplus_{p \mid d} A/\bigl(p^{v_p(d)}\bigr):
\qquad
T \cong \bigoplus_{p} \bigoplus_{j} A/\bigl(p^{k_{p,j}}\bigr),
\]
los \emph{divisores elementales}\index{divisores elementales}
$p^{k_{p,j}}$. Recíprocamente, los $d_i$ se reconstruyen a partir del
multiconjunto de \omterm{thm:b3:modules:structure}{divisores
elementales} ($d_s$ = producto de la mayor potencia de cada primo,
etc.), de modo que basta demostrar que el multiconjunto $\{k_{p,j}\}_j$
queda determinado por $T$ para cada primo $p$. Fijemos $p$; para
$j \geq 1$ considérense los $A/(p)$-espacios vectoriales
$p^{j-1}T/p^jT$. Sobre un factor cíclico $A/(p^k)$:
\[
p^{j-1}\bigl(A/(p^k)\bigr)\big/p^{j}\bigl(A/(p^k)\bigr)
\cong
\begin{cases}
A/(p) & \text{si } j \leq k,\\
0 & \text{si } j > k,
\end{cases}
\]
y sobre un factor $A/(q^k)$ con $q \neq p$: la multiplicación por $p$ es
allí biyectiva ($p$ es invertible \omterm{def:b3:modules:module}{módulo} $q^k$: Bézout), de modo que el
cociente es $0$. Las sumas directas se atraviesan: $\dim_{A/(p)}
p^{j-1}T/p^jT = \#\{i : k_{p,i} \geq j\}$. Estas dimensiones intrínsecas
determinan el multiconjunto de exponentes.
\end{proof}

\begin{corollary}[Grupos abelianos finitamente generados]
\label{cor:b3:modules:abelian}
Todo grupo abeliano \omterm{def:b3:modules:free}{finitamente generado}
es $\Z^r \times
\Z/d_1\Z\times\dots\times\Z/d_s\Z$ con $d_1 \mid \dots \mid
d_s$, de manera única. Todo grupo abeliano \emph{finito} es producto de
grupos cíclicos de orden potencia de primo, único como
multiconjunto.\index{grupo abeliano, estructura}
\end{corollary}

\begin{example}\label{ex:b3:modules:abelian}
Los grupos abelianos de orden $p^n$ corresponden a las
\emph{particiones} de $n$: para $p^4$: $\Z/p^4$, $\Z/p^3\times\Z/p$,
$(\Z/p^2)^2$, $\Z/p^2 \times (\Z/p)^2$, $(\Z/p)^4$ --- cinco grupos,
pues $4$ tiene cinco particiones. Los órdenes mixtos multiplican los
recuentos primo a primo (teorema chino del resto): hay $5 \times 2$
grupos abelianos de orden $2^4 \cdot 3^2 = 144$.
\end{example}

\section{Aplicación: formas canónicas de los endomorfismos}

Sea $K$ un cuerpo, sea $V$ un $K$-espacio vectorial de dimensión finita
$n$ y sea $u \in \mathcal L(V)$; hagamos de $V$ un $K[X]$-módulo
mediante $P
\cdot x = P(u)(x)$ (el~\cref{ex:b3:modules:examples}). Este
\omterm{def:b3:modules:module}{módulo} es
\omterm{def:b3:modules:free}{finitamente generado} (una $K$-base lo
genera) y de \omterm{def:b3:modules:torsion}{torsión}: para cada $x$,
los $n+1$ vectores $x, u(x), \dots, u^n(x)$ son $K$-dependientes, lo que
proporciona un polinomio anulador no nulo.

\begin{definition}\label{def:b3:modules:companion}
Para $P = X^m + a_{m-1}X^{m-1} + \dots + a_0$ mónico, la \emph{matriz
compañera}\index{matriz compañera} es
\[
C_P = \begin{pmatrix}
0 & & & -a_0\\
1 & \ddots & & -a_1\\
& \ddots & 0 & \vdots\\
& & 1 & -a_{m-1}
\end{pmatrix} :
\]
la matriz de «multiplicar por $X$» en $K[X]/(P)$ respecto de la base
$1, \bar X, \dots, \bar X^{m-1}$.
\end{definition}

\begin{theorem}[Frobenius: forma canónica racional]
\label{thm:b3:modules:frobenius}
Existe una única sucesión de polinomios mónicos no constantes $P_1
\mid P_2 \mid \dots \mid P_s$ (los \emph{invariantes de
semejanza}\index{invariantes de semejanza} de $u$) tal que, como
$K[X]$-módulos,
\[
V \cong K[X]/(P_1) \oplus \cdots \oplus K[X]/(P_s):
\]
y en una base adecuada $u$ tiene matriz diagonal por bloques
$\operatorname{diag}(C_{P_1}, \dots, C_{P_s})$. Además:
\begin{enumerate}
  \item $P_s = \mu_u$ (polinomio mínimo) y $P_1\cdots P_s =
        \chi_u$ (polinomio característico); en particular $\mu_u
        \mid \chi_u$ \emph{(nueva demostración de Cayley--Hamilton)} y
  $\chi_u \mid \mu_u^{\,s}$, de modo que $\chi_u$ y $\mu_u$ tienen los
  mismos factores \omterm{def:b3:rings:divisibility}{irreducibles}.
  \item Dos endomorfismos (o dos matrices cuadradas) son semejantes si y
  solo si tienen los mismos
  \omterm{thm:b3:modules:frobenius}{invariantes de semejanza}.
\end{enumerate}
\end{theorem}

\begin{proof}
El teorema de estructura (el~\cref{thm:b3:modules:structure}) aplicado
al \omterm{def:b3:rings:pidufd}{DIP} $K[X]$: el
\omterm{def:b3:modules:torsion}{módulo de torsión} $V$ se descompone con
\omterm{thm:b3:modules:smith}{factores invariantes} $P_i$, normalizados
mónicos (las unidades de $K[X]$ son $K^\times$); no aparece parte
\omterm{def:b3:modules:free}{libre} ($V$ es de
\omterm{def:b3:modules:torsion}{torsión}). Sobre cada factor cíclico
$K[X]/(P_i)$, la multiplicación por $X$ tiene matriz $C_{P_i}$ en la
base de las potencias de $\bar X$: concatenando bases se obtiene la
forma por bloques.

(1) El anulador de $V = \bigoplus K[X]/(P_i)$ es $(P_1)\cap
\dots\cap(P_s) = (P_s)$ (cadena de divisibilidad: $P_s$ es un múltiplo
común, y la clase de $1$ en el último factor es aniquilada exactamente
por $(P_s)$): $\mu_u = P_s$. Para $\chi_u$: sobre un factor cíclico,
$\chi_{C_P} = P$, por inducción sobre $m = \deg P$. Desarrollando
$\det(XI_m - C_P)$ por la primera fila (cuyas entradas son $X$, después
ceros y después $a_0$ en la última columna):
\[
\det(XI_m - C_P) = X\,\det\bigl(XI_{m-1} - C_{\tilde P}\bigr)
+ (-1)^{1+m}\,a_0\,\det L ,
\]
donde $\tilde P = X^{m-1} + a_{m-1}X^{m-2} + \dots + a_1$ (misma forma,
un tamaño menos) y $L$ es triangular con diagonal $(-1, \dots, -1)$,
luego $\det L = (-1)^{m-1}$. Por inducción, el primer término es
$X\tilde P$ y el segundo es $a_0$: el total es $X\tilde P + a_0 = P$
(caso base $m=1$: $\det(X + a_0) = P$). Los determinantes se multiplican
por bloques: $\chi_u = \prod P_i$. Cayley--Hamilton:
$\chi_u \in (\mu_u)$, pues $P_s \mid$ cada uno\dots\ recíprocamente,
cada $P_i \mid P_s$, luego $\chi_u = \prod P_i$ divide a
$P_s^{\,s} = \mu_u^s$; y $\mu_u =
P_s$ divide a $\chi_u$ por ser uno de sus factores.

(2) Dos endomorfismos semejantes son estructuras de
\omterm{def:b3:modules:module}{módulo} conjugadas, luego tienen los
mismos invariantes (unicidad en el~\cref{thm:b3:modules:structure});
recíprocamente, invariantes iguales dan $K[X]$-módulos isomorfos, y un
isomorfismo de \omterm{def:b3:modules:module}{módulos} es exactamente
una biyección lineal que entrelaza los dos endomorfismos: una semejanza.
\end{proof}

\begin{corollary}[La semejanza no cambia al extender el cuerpo]
\label{cor:b3:modules:descent}
Sean $K \subseteq L$ cuerpos y sean $M, N \in M_n(K)$. Si $M$ y $N$ son
semejantes sobre $L$, lo son sobre $K$.
\end{corollary}

\begin{proof}
Los \omterm{thm:b3:modules:frobenius}{invariantes de semejanza} de $M$
se calculan con la fórmula de los menores de Smith
(el~\cref{thm:b3:modules:smith}) aplicada a la matriz de presentación
$XI_n - M$ sobre $K[X]$ --- en efecto, el $K[X]$-módulo $V_M = K^n$
tiene presentación $XI_n - M$: la aplicación $K[X]^n \to V_M$,
$(Q_i) \mapsto \sum Q_i(M)e_i$, es sobreyectiva con núcleo generado por
las columnas de $XI_n - M$ (verificación directa: \omterm{def:b3:modules:module}{módulo} esas columnas,
todo elemento de $K[X]^n$ se reduce a un vector constante, y los
vectores constantes se aplican biyectivamente; el problema de fin de
semana lo detalla). Los máximos comunes divisores de polinomios no
cambian al extender el cuerpo: si $d$ es el máximo común divisor mónico
en $K[X]$ de una familia $(f_j)$, Bézout da $d = \sum u_jf_j$ con
$u_j \in K[X]$, de modo que todo divisor común de los $f_j$ \emph{en
$L[X]$} divide a $d$; y como $d$ es él mismo un divisor común, también
es el máximo común divisor en $L[X]$. Por tanto, los
\omterm{thm:b3:modules:smith}{factores invariantes} de $XI_n - M$,
cocientes de máximos comunes divisores de menores sucesivos, son los
mismos sobre $K$ y sobre $L$: $M, N$ tienen los mismos
\omterm{thm:b3:modules:frobenius}{invariantes de semejanza} sobre $L$ si
y solo si los tienen sobre $K$; se concluye por
\cref{thm:b3:modules:frobenius}(2).
\end{proof}

\begin{theorem}[Forma de Jordan,
redemostrada]\label{thm:b3:modules:jordan}
Supongamos que $\chi_u$ se descompone en factores lineales sobre $K$
(por ejemplo, $K = \C$). Aplicando a $V$ la descomposición en \omterm{thm:b3:modules:structure}{divisores
elementales} (demostración del~\cref{thm:b3:modules:structure}) en lugar
de la de \omterm{thm:b3:modules:smith}{factores invariantes}:
\[
V \cong \bigoplus_{\lambda, j} K[X]\big/\bigl((X -
\lambda)^{k_{\lambda,j}}\bigr),
\]
y en la base $\bigl(\overline{(X-\lambda)^{k-1}}, \dots,
\overline{(X - \lambda)}, \bar 1\bigr)$ de cada factor, $u$ actúa como
el bloque de Jordan $J_k(\lambda)$: todo endomorfismo cuyo polinomio
característico se descompone en factores lineales admite una base de
Jordan, y el multiconjunto de bloques $(\lambda, k)$ es
único.\index{forma de Jordan}
\end{theorem}

\begin{proof}
Los \omterm{thm:b3:modules:structure}{divisores elementales} del
\omterm{def:b3:modules:torsion}{módulo de torsión} $V$ son los
$(X - \lambda)^k$, donde $X - \lambda$ recorre los factores
\omterm{def:b3:rings:divisibility}{irreducibles} de $\mu_u$ (que se
descompone en factores lineales, pues así ocurre con $\chi_u$ y ambos
tienen los mismos factores
\omterm{def:b3:rings:divisibility}{irreducibles},
el~\cref{thm:b3:modules:frobenius}). En $W = K[X]/((X-\lambda)^k)$,
póngase $f_j = \overline{(X - \lambda)^{k-j}}$ para $j = 1, \dots, k$:
entonces $(X - \lambda)f_j = f_{j-1}$ (con $f_0 = 0$), es decir, $u(f_j)
= \lambda f_j + f_{j-1}$: la matriz de $u$ sobre $(f_1, \dots,
f_k)$ es exactamente $J_k(\lambda)$ (unos por encima de la diagonal). La
unicidad del multiconjunto de
\omterm{thm:b3:modules:structure}{divisores elementales} es
\cref{thm:b3:modules:structure}.
\end{proof}

\begin{remark}\label{rem:b3:modules:overview}
La jerarquía de las formas canónicas queda ahora clara: la forma
\emph{racional} existe sobre todo cuerpo y detecta la semejanza de
manera absoluta (el~\cref{cor:b3:modules:descent}); la forma de
\emph{Jordan} es su refinamiento cuando $\chi_u$ se descompone en
factores lineales. Las demostraciones del teorema de Jordan por recuento
de dimensiones vistas en segundo año quedan subsumidas: toda aquella
combinatoria era la aritmética del \omterm{def:b3:rings:pidufd}{DIP}
$K[X]$.
\end{remark}

\section{Ejercicios}

\begin{exercise}[$\star$]\label{exo:b3:modules:1}
(a) Demostrar que $\Q$ no es \omterm{def:b3:modules:free}{finitamente
generado} como $\Z$-módulo. (b) Demostrar que $\Q$ es
sin \omterm{def:b3:modules:torsion}{torsión} pero no
\omterm{def:b3:modules:free}{libre}. (c) ¿Por qué ninguno de los dos
enunciados contradice
\cref{thm:b3:modules:structure}?
\end{exercise}

\begin{exercise}[$\star$]\label{exo:b3:modules:2}
Listar los grupos abelianos de orden $360$ salvo isomorfismo, en forma
de \omterm{thm:b3:modules:structure}{divisores elementales} y en forma de \omterm{thm:b3:modules:smith}{factores invariantes}. ¿Cuántos
grupos abelianos de orden $p^5$ hay?
\end{exercise}

\begin{exercise}[$\star$]\label{exo:b3:modules:3}
Calcular la \omterm{thm:b3:modules:smith}{forma normal de Smith} sobre
$\Z$ de
\[
B = \begin{pmatrix} 2 & 4\\ 6 & 8 \end{pmatrix},
\qquad
C = \begin{pmatrix} 2 & 0 & 0\\ 0 & 3 & 0\\ 0 & 0 & 12
\end{pmatrix},
\]
e identificar los grupos abelianos $\Z^2/B\Z^2$ y $\Z^3/C\Z^3$.
\end{exercise}

\begin{exercise}[$\star\star$]\label{exo:b3:modules:4}
Sea $L \subseteq \Z^n$ un subgrupo de rango $n$ cuya base son las
columnas de $B \in M_n(\Z)$, $\det B \neq 0$. Demostrar que $\Z^n/L$ es
finito de cardinal $\abs{\det B}$, y que $\Z^n/L
\cong \prod_i \Z/d_i\Z$, donde $d_i$ son los
\omterm{thm:b3:modules:smith}{factores invariantes} de $B$. Ilústrese
con $L = \Z(2,0) + \Z(1,3)$.
\end{exercise}

\begin{exercise}[$\star\star$]\label{exo:b3:modules:5}
Sea $A$ un dominio. (a) Comprobar que $T(M)$ es un submódulo y que
$M/T(M)$ es sin \omterm{def:b3:modules:torsion}{torsión}. (b) Demostrar
que el \omterm{def:b3:rings:ideal}{ideal} $(X, Y)$ de $K[X,Y]$, visto
como $K[X,Y]$-módulo, es sin \omterm{def:b3:modules:torsion}{torsión}
pero no \omterm{def:b3:modules:free}{libre}: el teorema de estructura
necesita realmente la hipótesis de \omterm{def:b3:rings:pidufd}{DIP}.
\end{exercise}

\begin{exercise}[$\star\star$]\label{exo:b3:modules:6}
(a) Demostrar que $2\Z$ no tiene complemento directo en el $\Z$-módulo
$\Z$: los submódulos de los \omterm{def:b3:modules:free}{módulos libres}
son \omterm{def:b3:modules:free}{libres}
(el~\cref{thm:b3:modules:submodule}), pero no tienen por qué ser
sumandos directos. (b) Demostrar que si $M \subseteq A^n$ ($A$ un
\omterm{def:b3:rings:pidufd}{DIP}) cumple que $A^n/M$ es
sin \omterm{def:b3:modules:torsion}{torsión}, entonces $M$ \emph{sí} es
un sumando directo.
\end{exercise}

\begin{exercise}[$\star\star$]\label{exo:b3:modules:7}
Resolver en $\Z^2$ el sistema
\[
\begin{cases}
2x + 4y \equiv b_1 \pmod{20},\\
6x + 8y \equiv b_2 \pmod{20},
\end{cases}
\]
determinando para qué pares $(b_1, b_2)$ existen soluciones, mediante la
forma de Smith del~\cref{exo:b3:modules:3} (cambios de variable
invertibles en ambos lados).
\end{exercise}

\begin{exercise}[$\star\star$]\label{exo:b3:modules:8}
(a) Determinar todos los \omterm{thm:b3:modules:frobenius}{invariantes
de semejanza} y todas las \omterm{thm:b3:modules:jordan}{formas de
Jordan} posibles de una matriz nilpotente $4 \times 4$, clasificados
según la partición de $4$ que realizan. (b) Exhibir dos matrices
complejas $4\times4$ con el mismo polinomio característico \emph{y} el
mismo polinomio mínimo que no sean semejantes, y demostrar que para
$n \leq 3$ esto no puede ocurrir.
\end{exercise}

\begin{exercise}[$\star\star\star$]\label{exo:b3:modules:9}
Sean $u \in \mathcal L(V)$, $\dim V = n$. Demostrar que son
equivalentes: (i) $V$ es un $K[X]$-módulo cíclico (existe $x$ con
$V = K[u]x$, un \emph{vector cíclico}); (ii) $\mu_u = \chi_u$; (iii)
$s = 1$ en el~\cref{thm:b3:modules:frobenius}. Deducir que una
\omterm{def:b3:modules:companion}{matriz compañera} tiene un vector
cíclico, y determinar cuándo lo tiene una matriz diagonal.
\end{exercise}

\begin{exercise}[$\star\star\star$]\label{exo:b3:modules:10}
Para $M \in M_n(\Z)$, vista como endomorfismo de $\Z^n$, demostrar la
fórmula del índice: si $\det M \ne 0$, entonces
$[\Z^n : M\Z^n] = \abs{\det M}$, y deducir que $M \in GL_n(\Z)$ si y
solo si $\det M = \pm 1$. Aplicación: el grupo $\Z^2$ tiene exactamente
$\sigma_1(m) = \sum_{d \mid m} d$ subgrupos de índice $m$.
\emph{(Cuéntense las matrices en forma de Hermite
$\bigl(\begin{smallmatrix}
a & b\\ 0 & d\end{smallmatrix}\bigr)$, $ad = m$, $0 \leq b <
d$.)}
\end{exercise}

\begin{exercise}[$\star\star$]\label{exo:b3:modules:11}
(Ecuaciones $x^k = e$ en los grupos abelianos) Sea $G$ un grupo abeliano
finito con \omterm{thm:b3:modules:smith}{factores invariantes}
$d_1 \mid d_2 \mid \dots
\mid d_s$. (a) Demostrar que, para todo $k \geq 1$,
\[
\#\{x \in G : x^k = e\} \;=\; \prod_{i=1}^{s}\gcd(k, d_i) .
\]
(b) Deducir que un grupo abeliano finito es cíclico si y solo si, para
todo $k$, la ecuación $x^k = e$ tiene a lo sumo $k$ soluciones. (c)
Recuperar la ciclicidad de los subgrupos finitos de $K^\times$ ($K$ un
cuerpo, el~\cref{ch:b3:galois}): ¿por qué el polinomio $X^k - 1$
garantiza el criterio de (b)?
\end{exercise}

\begin{exercise}[$\star\star\star$]\label{exo:b3:modules:12}
(Las matrices elementales generan) (a) Demostrar que $M \in M_n(\Z)$ es
invertible en $M_n(\Z)$ si y solo si $\det M = \pm1$. (b) Demostrar que
$SL_2(\Z)$ está generado por las dos matrices elementales
$E = \bigl(\begin{smallmatrix}1 & 1\\ 0 &
1\end{smallmatrix}\bigr)$ y
$F = \bigl(\begin{smallmatrix}1 & 0\\ 1 &
1\end{smallmatrix}\bigr)$.
\emph{(Ejecútese el algoritmo de Euclides sobre la primera columna de $M
\in SL_2(\Z)$ mediante multiplicaciones por la izquierda por potencias
de $E, F$, hasta llegar a $\pm\bigl(\begin{smallmatrix}1 & *\\ 0 &
1\end{smallmatrix}\bigr)$; termínese a mano --- obsérvese que
$-I = (EF^{-1}E)^2$.)}
(c) Explicar la relación con la reducción de Smith: sobre $\Z$, las
operaciones de filas y columnas de determinante $1$ bastan para
diagonalizar, salvo signos.
\end{exercise}

\section{Problema: el conmutante y el biconmutante}

\begin{problem}[{Problema de fin de semana --- forma racional,
conmutante,
biconmutante}]\label{pb:b3:modules:1} Sea $K$ un cuerpo, sea $V$ un
$K$-espacio vectorial de dimensión $n \geq
1$ y sea $u \in \mathcal L(V)$. Estudiaremos el \emph{conmutante}
\[
\mathcal C(u) = \{v \in \mathcal L(V) : uv = vu\},
\]
una subálgebra de $\mathcal L(V)$ que contiene $K[u] = \{P(u) : P \in
K[X]\}$, y demostraremos la fórmula de dimensión de Frobenius y el
teorema del doble conmutante: $\mathcal C(\mathcal C(u)) = K[u]$. En
todo el problema, $V$ es el $K[X]$-módulo definido por $u$, con
\omterm{thm:b3:modules:smith}{factores invariantes}
$P_1 \mid \cdots \mid P_s$ y descomposición cíclica $V =
\bigoplus_{i=1}^s V_i$, $V_i = K[u]\,x_i \cong K[X]/(P_i)$, $n_i
= \deg P_i$ (el~\cref{thm:b3:modules:frobenius}).

\medskip
\noindent\textbf{Parte I --- La matriz de presentación $XI - M$ y
ejercicios preparatorios.}
\begin{enumerate}
  \item Sea $M \in M_n(K)$ y sea $\varphi \colon K[X]^n \to V_M
        = K^n$ la aplicación que envía $(Q_1, \dots, Q_n)$ a
  $\sum_i Q_i(M)e_i$. Demostrar que $\varphi$ es un morfismo
  sobreyectivo de $K[X]$-módulos y que toda columna de $XI_n - M$ está
  en $\ker\varphi$. \item Demostrar que, \omterm{def:b3:modules:module}{módulo} las columnas de
  $XI_n - M$, todo elemento de $K[X]^n$ es congruente con un vector
  constante \emph{(redúzcanse los grados usando
  $Xe_i \equiv \sum_j m_{ji}e_j$)}, y deducir
  $\ker\varphi = (XI_n - M)\,K[X]^n$: el
  \omterm{def:b3:modules:module}{módulo} $V_M$ tiene matriz de
  presentación $XI_n - M$. Recupérese el punto de partida
  del~\cref{cor:b3:modules:descent}: los
  \omterm{thm:b3:modules:frobenius}{invariantes de semejanza} de $M$ son
  los \omterm{thm:b3:modules:smith}{factores invariantes} no unidades de
  $XI_n - M$. \item Calcular los
  \omterm{thm:b3:modules:frobenius}{invariantes de semejanza} de: una
  matriz escalar $\lambda I_n$; una matriz diagonal con entradas
  diagonales distintas; el bloque de Jordan $J_n(0)$ de tamaño
  $n \times n$; $\operatorname{diag}(J_2(0), J_1(0))$ para $n = 3$.
  \item Demostrar que $\dim K[u] = \deg \mu_u = n_s$.
\end{enumerate}

\medskip
\noindent\textbf{Parte II --- Morfismos entre
\omterm{def:b3:modules:module}{módulos} cíclicos.}
\begin{enumerate}[resume]
  \item Sean $P, Q$ mónicos no constantes. Demostrar que un
  $K[X]$-morfismo $f \colon K[X]/(P) \to K[X]/(Q)$ queda determinado por
  $f(\bar 1)$, y que $\bar c \in K[X]/(Q)$ puede servir como $f(\bar 1)$
  si y solo si $P\bar c = 0$ en $K[X]/(Q)$. \item Deducir
        \[
        \operatorname{Hom}_{K[X]}\bigl(K[X]/(P),\, K[X]/(Q)\bigr)
        \;\cong\; K[X]\big/\bigl(\gcd(P, Q)\bigr),
        \]
        de dimensión $\deg \gcd(P, Q)$ sobre $K$. \emph{(Demuéstrese que
        las soluciones $\bar c$ de $P\bar c = 0$ en $K[X]/(Q)$ forman el
        submódulo cíclico generado por
        $\overline{Q/\gcd(P,Q)}$.)}
  \item Demostrar la \emph{fórmula de Frobenius}:
        \[
        \dim_K \mathcal C(u) = \sum_{i,j=1}^{s} \deg\gcd(P_i,
        P_j) = \sum_{i=1}^{s} (2s - 2i + 1)\, n_i .
        \]
        \emph{(Una $v$ que conmuta es exactamente un $K[X]$-endomorfismo
        de $V$; descompóngase $\operatorname{End}(\bigoplus_i V_i)$ en
        matrices de morfismos $V_j \to V_i$ y úsese la cadena de
        divisibilidad.)} \item Deducir $\dim \mathcal C(u) \geq n$, con
        igualdad si y solo si $u$ es cíclico ($s = 1$), y calcular
        $\dim\mathcal C(u)$ para $u = \lambda\,\mathrm{id}$: los dos
        extremos de la fórmula. \item Verificar directamente la fórmula
        de Frobenius para $\operatorname{diag}(J_2(0), J_1(0))$,
        calculando el conmutante explícitamente como matrices
        $3\times3$.
\end{enumerate}

\medskip
\noindent\textbf{Parte III --- El teorema del doble conmutante.} Sea
$w \in \mathcal C(\mathcal C(u))$; demostraremos que $w \in K[u]$.
\begin{enumerate}[resume]
  \item Demostrar que $K[u] \subseteq \mathcal C(\mathcal C(u))$ y que
  todo $w \in \mathcal C(\mathcal C(u))$ conmuta con $u$ --- de modo que
  la inclusión por demostrar,
  $\mathcal C(\mathcal C(u)) \subseteq K[u]$, es un refinamiento genuino
  de $w \in \mathcal C(u)$. \item Supongamos primero que $u$ es
  \emph{cíclico}, $V = K[u]x$. Demostrar directamente que
  $\mathcal C(u) = K[u]$ \emph{(evalúese una $v$ que conmute en $x$:
  $v(x) = P(u)x$ para cierto $P$, y compárense $v$ y $P(u)$ sobre la
  base $u^k x$)}, y concluir el teorema en este caso. \item Volvamos al
  caso general. Para cada $i$, sea $\pi_i\colon
        V \to V_i$ la proyección a lo largo de los demás sumandos.
  Demostrar que $\pi_i \in \mathcal C(u)$ y deducir que $w$ conserva
  cada $V_i$ y conmuta con $u_i =
        u\restriction_{V_i}$; concluir mediante la pregunta 11 aplicada
  al $u_i$ cíclico: existen polinomios $Q_i$ con
        $w\restriction_{V_i} = Q_i(u)\restriction_{V_i}$.
  \item Queda pegar los $Q_i$ en un único polinomio. Para $i
        \leq j$ (de modo que $P_i \mid P_j$), demostrar que
  $\eta_{ij} \colon
        V_j \to V_i$, $R(u)x_j \mapsto R(u)x_i$, es un $K[X]$-morfismo
  bien definido \emph{(lo que hay que comprobar es que $R(u)x_j = 0$
  implica $R(u)x_i = 0$)}, y que
  $\tilde\eta_{ij} = \eta_{ij}\circ\pi_j$, extendida por $0$ sobre los
  demás sumandos, está en $\mathcal C(u)$. \item Usando
  $w\tilde\eta_{ij} = \tilde\eta_{ij}w$, demostrar que $Q_i
        \equiv Q_j \pmod{P_i}$ para $i \leq j$. Deducir que $Q =
        Q_s$ cumple $Q \equiv Q_i \pmod {P_i}$ para todo $i$, luego
  $w = Q(u)$ sobre cada $V_i$ y, por tanto, sobre $V$:
        \[
        \boxed{\ \mathcal C(\mathcal C(u)) = K[u].\ }
        \]
  \item (Coda) Deducir del teorema: si $v$ conmuta con \emph{toda}
  matriz que conmute con $u$ y $u$ es cíclico, entonces $v$ es un
  polinomio en $u$; y dar un ejemplo que muestre que
  $\mathcal C(u) = K[u]$ falla para $u =
        \mathrm{id}$, $n \geq 2$ --- ¿dónde interviene exactamente la
  ciclicidad?
\end{enumerate}

\medskip
\noindent\textbf{Parte IV --- Dividendos de los
\omterm{thm:b3:modules:frobenius}{invariantes de semejanza}.} La forma
canónica racional es una máquina; he aquí cinco de sus resultados
clásicos.
\begin{enumerate}[resume]
  \item (Traspuesta) Demostrar que toda $M \in M_n(K)$ es semejante a su
  traspuesta ${}^tM$. \emph{(Las operaciones que llevan $XI - M$ a forma
  de Smith, traspuestas, llevan $XI - {}^tM$ a la misma forma de Smith:
  mismos \omterm{thm:b3:modules:frobenius}{invariantes de semejanza}.)}
  \item (Descenso de la semejanza) Sea $K \subseteq L$ una extensión de
  cuerpos y sean $M, N \in M_n(K)$. Demostrar que si $M$ y $N$ son
  semejantes sobre $L$, lo son sobre $K$. \emph{(La forma de Smith de
  $XI - M$ calculada en $K[X]$ sigue siendo una forma de Smith en $L[X]$
  --- ¿por qué no cambian los \omterm{thm:b3:modules:smith}{factores
  invariantes}?)} Consecuencia que conviene recordar: dos matrices
  reales conjugadas en $GL_n(\C)$ lo son en
        $GL_n(\R)$.
  \item (Clasificación de los nilpotentes) Sea $u$ nilpotente. Demostrar
  que el número de bloques de tamaño $\geq k$ en su descomposición en
  bloques de Jordan nilpotentes es igual a
  $\operatorname{rk}u^{k-1} - \operatorname{rk}u^k$, y deducir que las
  clases de nilpotentes de $M_n(K)$, sobre cualquier cuerpo $K$, están
  en biyección con las particiones de $n$. ¿Cuántas clases de
  nilpotentes hay en $M_5(K)$? \item (Un par concreto) Determinar los
  \omterm{thm:b3:modules:frobenius}{invariantes de semejanza} de la
  derivación $D\colon P \mapsto P'$ actuando sobre el espacio
  $K_{n-1}[X]$ de los polinomios de grado $< n$: (a) para $K =
        \Q$; (b) para $K = \mathbb F_p$ con $p < n$ \emph{(en
  característica $p$, $(X^p)' = 0$: calcúlese $\ker D^k$ y úsese la
  pregunta 18)}. \item (\omterm{ex:b3:groups:actions}{Clases de
  conjugación} de $GL_2(\mathbb F_q)$) Usando los
  \omterm{thm:b3:modules:smith}{factores invariantes}, demostrar que
  toda clase de $GL_2(\mathbb F_q)$ es exactamente de uno de estos
  cuatro tipos: central $aI$; diagonalizable con dos valores propios
  distintos $a \neq b$ en $\mathbb F_q^\times$; no semisimple con
  polinomio mínimo $(X - a)^2$; cíclica con polinomio característico
  \omterm{def:b3:rings:divisibility}{irreducible}. \item Contar las
  clases de cada tipo y concluir: $GL_2(\mathbb F_q)$ tiene exactamente
  $q^2 - 1$ \omterm{ex:b3:groups:actions}{clases de conjugación}.
  \emph{(Cuéntense las cuádricas mónicas
  \omterm{def:b3:rings:divisibility}{irreducibles} sobre $\mathbb F_q$;
  los pares no ordenados $\{a, b\}$; recuérdese que la invertibilidad
  restringe los términos independientes.)} \item (Lo cíclico es
  genérico) Demostrar que $M \in M_2(\mathbb F_q)$ deja de ser cíclico
  si y solo si $M$ es escalar, y deducir que una matriz $2\times2$
  tomada al azar uniformemente sobre $\mathbb F_q$ es cíclica con
  probabilidad $1 - q^{-3}$. Enunciar la heurística análoga para $M_n$ y
  $q$ grande (no se pide demostración): las matrices no cíclicas son
  raras --- por eso la Parte III del~\cref{pb:b3:modules:1} solo
  requirió trabajo de verdad más allá del caso genérico.
\end{enumerate}

\medskip
\noindent\textbf{Parte V --- Complementos.}
\begin{enumerate}[resume]
  \item (\omterm{ex:b3:groups:actions}{Centro} del conmutante) Demostrar
  que el \omterm{ex:b3:groups:actions}{centro} del álgebra
  $\mathcal C(u)$ es exactamente $K[u]$ \emph{(combínense las dos
  inclusiones de la Parte III)}. Deducir que $\mathcal C(u)$ es
  conmutativa si y solo si $u$ es cíclico --- lo que recupera el caso de
  igualdad de la pregunta 8 por una vía puramente estructural. \item
  (¿Qué dimensiones se alcanzan?) Deducir de la fórmula de Frobenius que
  $\dim\mathcal C(u) \equiv n \pmod 2$ para todo $u$. Determinar después
  el conjunto exacto de valores que toma $\dim \mathcal C(u)$ cuando $u$
  recorre $\mathcal L(V)$ con $\dim V = 4$: demostrar que es $\{4, 6,
        8, 10, 16\}$ \emph{(enumérense las sucesiones de grados $n_1
        \leq \dots \leq n_s$ que suman $4$ y realícese cada una con un
  nilpotente)}. En particular, $12$ y $14$, pese a tener la paridad
  correcta, no se alcanzan: la restricción de paridad es necesaria, pero
  no suficiente. \item (\omterm{cor:b3:groups:classeq}{Ecuación de
  clases} de $GL_2(\mathbb F_3)$) Para $q = 3$, calcular el tamaño de
  cada \omterm{ex:b3:groups:actions}{clase de conjugación} de la
  pregunta 20 mediante órbita--estabilizador: el centralizador de una
  $M$ \emph{cíclica} en $GL_2(\mathbb F_q)$ es el grupo de unidades de
  $K[M]$ (pregunta 11). Identificar $K[M]$ en los tres tipos no
  centrales, listar los tres polinomios cuadráticos mónicos
  \omterm{def:b3:rings:divisibility}{irreducibles} sobre $\mathbb F_3$ y
  verificar la \omterm{cor:b3:groups:classeq}{ecuación de clases}
        \[
        48 = \abs{GL_2(\mathbb F_3)}
        = 2\cdot1 + 1\cdot12 + 2\cdot8 + 3\cdot6,
        \]
        con $2 + 1 + 2 + 3 = 8 = q^2 - 1$ clases, tal como predice la
        pregunta 21.
\end{enumerate}
\end{problem}
